\section{Flujos principales}

\subsection{Matrícula nueva y renovación}

Buscar por nombre o RUN/IPE/IPA $\rightarrow$ Nueva o Renovar $\rightarrow$
wizard de 8 pasos $\rightarrow$ documentos $\rightarrow$ firma $\rightarrow$
confirmación. La renovación precarga estudiante, familia, apoderados,
contactos, retiro, salud, PIE y autorizaciones; si el estudiante terminó
8º básico pide confirmar egreso o repetición.

Bloques del wizard: A estudiante, B familia y apoderados, C emergencia y
retiro, D salud, E PIE y apoyos, F autorizaciones de imagen, G convivencia,
paso 8 revisión con PDF embebido y confirmación expresa del funcionario.

Estados: \codigo{BORRADOR} $\rightarrow$ \codigo{FICHA\_GENERADA} $\rightarrow$
\codigo{FIRMADA} $\rightarrow$ \codigo{CONFIRMADA} (o \codigo{ANULADA} con
motivo). La base impide saltos; la anulación permite iniciar una matrícula
nueva. La columna \codigo{version} actúa como token de concurrencia: si otro
funcionario guardó primero, EF lanza excepción en vez de sobrescribir.

\subsection{Ficha y PDF}

Al cerrar se crea \codigo{ficha\_version} con folio, versión y
\codigo{datos\_snapshot}: el PDF se genera desde la copia congelada, nunca
desde el borrador. La corrección crea una versión nueva con motivo y deja la
anterior \codigo{REEMPLAZADA} sin tocar el documento firmado. La descarga
queda registrada como \codigo{DESCARGAR}.

\subsection{Expediente documental}

Carga PDF/JPG/PNG hasta 15 MB con validación de extensión y firma real,
hash SHA-256 y almacenamiento fuera de \codigo{wwwroot}. Cada requisito se
marca recibido/no aplica/pendiente sin bloquear la matrícula. El escaneo de
la ficha firmada se guarda como documento \codigo{FICHA\_FIRMADA};
\codigo{AlumnosController.ValidarDocumentacion} registra ubicación del
original, validador y observación.

\subsection{Autorizaciones, certificados y alertas}

Las autorizaciones conservan firma y fecha; la revocación exige motivo y
bloquea el uso posterior, y las vencidas pasan a \codigo{VENCIDA}. El
certificado de alumno regular se emite desde los datos confirmados, con
folio y código, historial y validación pública por folio más código
(\codigo{/c/\{código\}}); la anulación le quita validez. \codigo{/Alertas} lista
vencidas y las que vencen en 30 días, filtradas por año y por área
(\codigo{AUTORIZACIONES:LEER} / \codigo{MEDICACION:LEER}); la revisión
programada corre cada 15 minutos. Los reportes (\codigo{ReportesController})
exportan CSV de matrículas y faltantes con permiso EXPORTAR y descarga
auditada.

\subsection{Cuentas}

ADMIN\_TI crea, activa, desactiva, cierra rol, restablece contraseña (clave
temporal de un solo uso con cierre de sesiones) y desbloquea cuentas desde
\codigo{FuncionariosController}.

\clearpage
\subsection{Diagramas de los procesos}

\diagramaAncho{03-proceso-login}{1.0}{Inicio de sesión: cookie, permiso y sesión revocada.}
\diagramaAncho{04-proceso-matricula}{1.0}{Matrícula: ficha, folio, firma, confirmación y validación.}
\diagramaAncho{05-proceso-certificado}{1.0}{Certificado: emisión, folio y validación pública.}
\diagramaAncho{06-proceso-medicacion}{1.0}{Medicación: registro de dosis y triggers.}

\clearpage
